\begin{mistake}{2026-05-27}{04}

\begin{problemcontent}
设 \(D\) 为圆域 \(x^2 + y^2 \leqslant 2x + 2y\)，则 \(\iint_D xy \,\mathrm{d}x\mathrm{d}y = \underline{\hspace{2cm}}\).
\end{problemcontent}

\begin{solutioncontent}
首先对积分区域进行配方：
\[(x-1)^2 + (y-1)^2 \leqslant 2\]
区域 \(D\) 为圆心在 \((1,1)\)，半径为 \(\sqrt{2}\) 的圆。

作平移变换，令 \(u = x-1\), \(v = y-1\)，则 \(x = u+1\), \(y = v+1\)。
积分区域化为以原点为中心的圆域 \(D' : u^2 + v^2 \leqslant 2\)。

代入原积分并展开被积函数：
\[\iint_D xy \,\mathrm{d}x\mathrm{d}y = \iint_{D'} (u+1)(v+1) \,\mathrm{d}u\mathrm{d}v = \iint_{D'} (uv + u + v + 1) \,\mathrm{d}u\mathrm{d}v\]

由于 \(D'\) 关于 \(u\) 轴和 \(v\) 轴均对称，且被积函数中的 \(uv\)、\(u\)、\(v\) 在对称区域上对应的积分为 \(0\)。
原积分仅剩常数项 \(1\) 的积分：
\[\iint_{D'} 1 \,\mathrm{d}u\mathrm{d}v = S_{D'} = \pi (\sqrt{2})^2 = 2\pi\]
故答案为 \(2\pi\)。
\end{solutioncontent}

\mistakeknowledge{
\begin{itemize}
  \item \textbf{核心公式/定理}：二重积分平移变换定理、积分域对称性与奇偶函数性质。
  \item \textbf{易错点}：未利用平移变换而直接采用极坐标计算导致运算量极大；平移后展开 \((u+1)(v+1)\) 时遗漏交叉项。
  \item \textbf{关联知识}：被积函数为 \(1\) 的二重积分在数值上等于积分区域的面积。
\end{itemize}}

\end{mistake}
